% Préambule commun aux cours IPR Scholar.
\documentclass[11pt,a4paper,oneside]{article}
\usepackage[T1]{fontenc}
\usepackage[utf8]{inputenc}
\usepackage[french]{babel}
\usepackage{lmodern}
\usepackage{microtype}
\usepackage[a4paper,margin=2.35cm]{geometry}
\usepackage{amsmath,amssymb,mathtools}
\usepackage{booktabs,longtable,array}
\usepackage{xcolor}
\usepackage{hyperref}
\usepackage{tikz}
\usetikzlibrary{arrows.meta,positioning,shapes.geometric,calc}
\hypersetup{colorlinks=true,linkcolor=blue!55!black,citecolor=blue!55!black,urlcolor=blue!55!black}
\definecolor{iprblue}{HTML}{1F4E79}
\definecolor{iprgreen}{HTML}{2D6A4F}
\newcommand{\IPR}{\textsc{ipr}}
\newcommand{\coursmeta}[2]{\author{Projet \IPR{} -- documentation Scholar}\date{#1}\maketitle\begin{center}\small #2\end{center}\vspace{0.5cm}\tableofcontents\clearpage}
\newcommand{\resultat}[1]{\paragraph{À retenir.} #1}

\title{Cours 8 -- Signal, bruit et grain déterministe}
\begin{document}
\coursmeta{5 octobre 2026}{Synthèse de grain et débruitage CPU intégré à \IPR{} ; validation photographique du correctif à refaire.}

\section{Théorie du bruit photographique}
Le bruit photonique suit approximativement une loi de Poisson : pour $N$ photons, l'écart-type est $\sqrt{N}$ et le rapport signal/bruit vaut $\sqrt{N}$. S'ajoutent le bruit de lecture, le bruit thermique et les erreurs de quantification. Le grain argentique et le bruit numérique ne sont donc pas identiques : le premier a une structure spatiale et chromatique dépendante de l'émulsion, le second dépend du capteur et du dématriçage~\cite{gonzalezWoods}.

\section{Application dans \IPR{}}
Le réglage Grain ne simule pas un capteur. Il synthétise une perturbation monochrome reproductible : une fonction de hachage des coordonnées $(x,y)$ donne un bruit uniforme $n\in[-1,1]$. La luminance devient
\begin{equation}
Y'=\max(0,Y+n\,0{,}06\,g\,w_s(Y)),
\end{equation}
où $g\in[0,1]$ est le curseur et $w_s(Y)=0{,}45+0{,}55(1-\operatorname{clamp}(Y,0,1))$. Les ombres reçoivent donc davantage de grain.

\section{Pourquoi un hash déterministe}
L'aperçu peut être recalculé plusieurs fois par seconde. Un générateur aléatoire non déterministe ferait scintiller le grain et rendrait impossible une comparaison rigoureuse aperçu/export. Les constantes entières de mélange servent à répartir les bits; elles n'ont pas de signification photographique. La valeur $0{,}06$ est un plafond esthétique empirique : elle empêche le réglage maximal de masquer la scène.

\section{Exemple}
À $Y=0{,}10$ et Grain 50, l'amplitude maximale est environ $0{,}06\times0{,}5\times0{,}945\simeq0{,}028$. À $Y=0{,}90$, elle descend vers $0{,}016$. Le grain est ainsi plus lisible dans un mur sombre que dans un ciel clair.

\section{Évolution possible}
Un grain plus crédible demande une corrélation spatiale, une taille dépendante de la résolution et éventuellement une composante chromatique. La réduction de bruit est une tâche différente : un module CPU est désormais intégré, mais sa qualité photographique ne se déduit pas de la seule baisse d'une statistique de bruit.

\section{Débruitage : estimateur et seuillage}
Les plans du moteur courant relu le 22 octobre 2026 sont Y Rec.2020 linéaire, $Co=R-B$ et $Cg=G-(R+B)/2$.
Le coefficient diagonal Haar
\begin{equation}
h=\frac{p_{00}-p_{01}-p_{10}+p_{11}}{2}
\end{equation}
a une norme quadratique unité : pour quatre bruits gaussiens indépendants
de variance $\sigma^2$, sa variance reste $\sigma^2$. L'estimateur
$\widehat\sigma=\operatorname{median}(|h-\operatorname{median}(h)|)/0{,}67448975$
utilise la médiane de la valeur absolue d'une normale centrée réduite.
Après dématriçage, le bruit est corrélé : l'interprétation physique exacte
n'est plus garantie. L'implémentation échantillonne des blocs adjacents et
ignore les blocs parfaitement constants, sans prétendre détecter les pixels chauds.

Deux flous successifs séparent les résidus fin et moyen. Le seuillage doux
$S_\lambda(r)=\operatorname{sgn}(r)\max(|r|-\lambda,0)$ est piloté par
$\lambda=\widehat\sigma ms$, où $s$ est la valeur divisée par 100 et
$m$ vaut 2,55 en luminance, 3,15 en chrominance. Le seuil moyen vaut
0,52 fois le seuil fin. Les seuils vont à zéro à l'activation ; le mélange
final avec l'original, de poids $s$, garde les états intermédiaires utiles.

\section{Structure, luminosité et exemple}
Le guide de contours est lissé avant son Sobel, pour éviter que le bruit
soit protégé comme une texture. Les changements signés sont rééquilibrés
pour conserver la moyenne ; le gamut négatif est comprimé vers un gris
de même luminance, pas coupé indépendamment par canal. Le plancher
historique 0,00025 reste uniquement dans la netteté, pas dans le débruitage.

Dans l'essai synthétique historique du 5 octobre, sur un aplat de signal 0,05 et bruit gaussien monochrome
$\sigma=0{,}005$, le proxy Haar restant est environ 0,831, 0,523, 0,253
et 0,009 fois celui du neutre pour les réglages 25/50/75/100. Cet exemple
ne contient pas de texture à préserver : une baisse aussi forte ne valide
ni un portrait ni un feuillage. Une moyenne globale stable ne garantit pas
la stabilité des régions sombres. Les nouvelles photographies doivent
valider précisément ces limites.

\resultat{Le grain actuel est un effet créatif reproductible, pas un modèle de bruit de capteur ni un algorithme de débruitage.}
\bibliographystyle{plain}\bibliography{references}
\end{document}
